\section{Quick start}

\subsection{The GUI}

Double-click \cmd{Run Simulation GUI.bat} in the repository root. The
first launch installs a private Python runtime into \cmd{gui/runtime/}
(internet required, a few minutes) with the repository's pinned
package versions plus the GUI extras; afterwards it starts offline,
creates an icon-bearing shortcut, and opens at
\cmd{http://localhost:8552}.

The studio walks the pipeline as a gated workflow: Specimen \&
Fabric, Probe \& Excitation, Error model, then Acquisition; the
Backscatter and COF reconstruction tabs unlock once acquired data
exists. The FW sweep tab drives the resumable full-wave azimuth sweep
engine.

Everything interactive (specimen builder, ray forward model, error
model, reconstruction views) runs on CPU. New full-wave runs need
CUDA and cupy in the running Python environment; the GUI detects the
backend by inspection and reports it under the title.

\subsection{Reproducing a stored result}

No GPU is needed to reproduce any analysis in the paper. For example,
the finite-element cross-check of direct-arrival attenuation:

\begin{lstlisting}
cd sim/analysis
python fe_direct_attenuation.py
\end{lstlisting}

prints a table of the five cross-check rounds; the expected values,
including the 0.027 dB solver agreement for the 6 to 15 mm placement,
are recorded in the module docstring. Every cited module follows the
same pattern: run it from its own directory, compare the printed
numbers with the docstring.
